\documentclass[11pt]{article}
\usepackage[T1]{fontenc}
\usepackage{lmodern,amsmath,amssymb,amsthm,microtype}
\usepackage[margin=1in]{geometry}
\numberwithin{equation}{section}
\usepackage[hidelinks]{hyperref}
\newtheorem{theorem}{Theorem}[section]
\newtheorem{corollary}[theorem]{Corollary}
\newtheorem{lemma}[theorem]{Lemma}
\newtheorem{definition}[theorem]{Definition}
\newcommand{\E}{\mathbb E}
\newcommand{\R}{\mathbb R}
\newcommand{\one}{\mathbf 1}
\newcommand{\ip}[2]{\langle #1,#2\rangle}
\newcommand{\Ga}{\mathcal R_\alpha}
\newcommand{\Cspin}{\mathcal C_n}
\allowdisplaybreaks
\setlength\emergencystretch{2em}
\title{Gaussian and non-Gaussian finite-size scalings of the
spread-out Ising model at the free-boundary pseudocritical point}
\author{Jiwoon Park}
\date{\today}
\begin{document}
\maketitle

\begin{abstract}
For the critical weakly coupled $|\varphi|^4$ model with periodic
boundary conditions, Park \cite{park-torus} showed that only the
constant mode has enhanced fluctuations. Motivated by the shifted
free-boundary critical window established by Michta, Park, and Slade
\cite{michta-park-slade} for the hierarchical model, we study the
spread-out Ising model at its free-boundary pseudocritical point.
At the upper critical dimension we determine the exponential moments
of the spin field tested against smooth real-valued functions vanishing
on the boundary. For a nonzero test, the spin sum divided by its
standard deviation has a quartic limiting law if the test overlaps
the ground state and a standard Gaussian limit if it is orthogonal.
In four dimensions the orthogonal susceptibility has order $n^2$
in a box of side $n$, the scale obtained by Camia, Jiang, and Newman
\cite{camia-jiang-newman} at bulk criticality under a mean-field
two-point hypothesis. The lowest sine and the $(3,1,1,1)$ sine have
different limiting laws, so it provides a counterevidence to the all-odd 
classification of Flores-Sola, Berche, Kenna, and Weigel \cite{flores}.
\end{abstract}


\subsubsection*{Version note: you can find earlier versions on DOI:10.5281/zenodo.24167641}

v1. November 20, 2026

\setcounter{tocdepth}{1}
\tableofcontents

\section{The model and the two susceptibility scales}

We consider the spread-out Ising model on a box with free boundary condition:
fix $d\in\{1,2,3,4\}$, $\alpha=d/4$, $s=2\alpha=d/2$,
$\Lambda_n=\{1,\ldots,n\}^d$ and $V=n^d$.
Let $D_R$ be the probability measure uniform on $0<|z|_\infty\le R$ in
$\mathbb Z^d$, and set $L_R=I-D_R$.
For inverse temperature $\tau >0$, we consider the measure on $\sigma \in\{-1,1\}^{\Lambda_n}$
with weight proportional to
\begin{align}
 \exp\{-\tfrac12\tau\,\sigma^{\mathsf T}J_n\sigma\},\qquad
 J_n=P_{\Lambda_n} L_R^\alpha P_{\Lambda_n}^* .
 \label{eq:Hamiltonian}
\end{align}
where $P_{\Lambda_n} : \ell^2 (\mathbb{Z}^d) \rightarrow \ell^2 (\Lambda_n)$ 
is the simple restriction to $\Lambda_n$ and $P_{\Lambda_n}^*$ extends by zero, 
so the fractional power is taken before restriction to the box. 
Fix $R$ sufficiently large. 
For $n\ge n_0(R)$ choose the pseudocritical point
$\tau=\tau^{\rm pc}_{\Lambda_n}$ constructed in \cite{base}. At this
temperature the negative logarithm of the marginal density of the
prescribed terminal coordinate has zero curvature at the origin.
Section~2 of \cite{base} defines its principal direction.
Flores-Sola et al.\ describe their pseudocritical temperature $T_L$
using a finite-volume response peak \cite{flores}.

At the free-boundary pseudocritical point, the variance of the total
spin divided by the volume is of order $n^{d/2}\sqrt{\log n}$ in a box
of side $n$ \cite{base}. 
Does the second moment carry the factor $\sqrt{\log n}$ for any test function?
The question comes from Table~I of Flores-Sola, Berche, Kenna and
Weigel \cite{flores}.
% , which assigns different pseudocritical scaling sector for $d>d_c$ according to the parity of Fourier mode of the test function.
At $d=d_c$ we find that the logarithmic enhancement is determined by overlap with the ground state. 
In the short-range model this singles out the lowest sine. In the long-range
models the ground state is a superposition of infinitely many sines.
The difference between overlap with the uniform field and overlap with
the ground state can be seen explicitly in \eqref{eq:odd-pair}.

Unless an inverse temperature is displayed explicitly, spin expectations
are taken at the pseudocritical point.

\begin{definition}[Ising measure at $\tau_{\Lambda_n}^{\mathrm{pc}}$]
Write $\E_{\rm sp}$ for expectation under this spin law given by the weight \eqref{eq:Hamiltonian}
and $G_{\rm sp}(x,y)=\E_{\rm sp}[\sigma_x\sigma_y]$ for its two-point function
having \emph{\textbf{fixed} $\tau = \tau_{\Lambda_n}^{\mathrm{pc}}$}.
\end{definition}

We start with a remark that $\tau_{\Lambda_n}^{\mathrm{pc}}$ is a genuine shift from the
bulk cricial point 
$\tau_c=\tau_c(d,R)=\inf\{\tau>0:\langle\sigma_0\rangle_\tau^+>0\}$,
where $\langle\cdot\rangle_\tau^+$ is the infinite-volume plus state for
the interaction in \eqref{eq:Hamiltonian}. Write
$\Delta_n=\tau_{\Lambda_n}^{\rm pc}-\tau_c$ and $(u)_+=\max\{u,0\}$.

\begin{theorem}[Shift of the pseudocritical point]\label{thm:pc-shift}
For $d\in\{1,2,3,4\}$ and sufficiently large fixed $R$,
$\tau_{\Lambda_n}^{\rm pc}\to\tau_c$ as $n\to\infty$.
For $n\ge n_1(d,R)$ one has $(\Delta_n)_+\le C_dR^{-d}$.
For $d=1,2,3$ one also has
\begin{equation}\label{eq:pc-shift-upper}
 (\Delta_n)_+\le C_{d,R}
 \begin{cases}
   n^{-1/10}(\log n)^{1/10}, & d=1,3,\\
   n^{-1/5}(\log n)^{1/2}, & d=2 .
 \end{cases}
\end{equation}
There is $c_{d,R}>0$ such that, for all sufficiently large $n$,
\begin{equation}\label{eq:pc-shift}
 \Delta_n\ge c_{d,R}\,n^{-s}=c_{d,R}\,n^{-d/2}.
\end{equation}
\end{theorem}

Refinement of this theorem is of seperate interest. Since the theorem is
independent of the rest of the article, we isolate the proof in Appendix~\ref{app:shift}.


\paragraph{Spin observables.}

All lattice and continuum test functions are real-valued. For lattice
vectors we use the counting-measure scalar product
$\ip{f}{h}=f^{\mathsf T}h$. The spin sum is
$S_n(f)=\ip{f}{\sigma}=\sum_{x\in\Lambda_n} f(x)\sigma_x$.
On $\Omega = (0,1)^d$, define the continuum sines for
$k\in\mathbb N^d$, where $\mathbb N=\{1,2,\ldots\}$, and their
lattice samples for $1\le k_i\le n$ by
\begin{equation}
 e_k(z)=2^{d/2}\prod_{i=1}^d\sin(\pi k_i z_i),\qquad
 f_n(F)(x)=(n+1)^{-d/2}F(x/(n+1)),\qquad \phi_{k,n}=f_n(e_k).
\end{equation}
The lattice sines are orthonormal and vanish on the Dirichlet planes
$x_i=0,n+1$. Write $k_g=(1,\ldots,1)$, $e_g=e_{k_g}$ and
$\phi_{g,n}=\phi_{k_g,n}$. The covariance and $F$-susceptibility of
sampled tests are
\begin{equation}
 \Cspin(F,H)=\E_{\rm sp}[S_n(f_n(F))S_n(f_n(H))]
             =f_n(F)^{\mathsf T}G_{\rm sp} f_n(H),\qquad \chi_n(F)=\Cspin(F,F).
\end{equation}
Thus $\chi_{k,n}=\chi_n(e_k)$ is the susceptibility of a unit sine,
and the lattice test $V^{-1/2}\one$ gives the uniform susceptibility.

\paragraph{The continuum operator.}
Let $\Omega = (0,1)^d$ be equipped with Lebesgue measure.
Let $H_\alpha$ be the self-adjoint operator of the closed form
\begin{equation}
 \mathcal E_\alpha(F,F)=\int_{\R^d}|\xi|^{2\alpha}
       |\widehat{\widetilde F}(\xi)|^2\,d\xi ,
 \qquad \widetilde F\in H^\alpha(\R^d),\quad
 \widetilde F=0 \text{ off }\overline\Omega ,
 \label{eq:form}
\end{equation}
where $\widetilde F$ is the zero extension of $F$ and
$\widehat u(\xi)=(2\pi)^{-d/2}\int e^{-ix\cdot\xi}u(x)\,dx$.
Write $\mu_1<\mu_2\le\cdots$ for the eigenvalues, counted with
multiplicity, and $w$ for the nonnegative $L^2$-unit ground state.
The reduced resolvent $\Ga$ inverts $H_\alpha-\mu_1$ on $w^\perp$.
We use
\begin{equation}
 \Ga w=0,\qquad
 \Ga|_{w^\perp}=(H_\alpha-\mu_1)^{-1}|_{w^\perp},\qquad
 I_4=\int_\Omega w^4,\qquad b_F=\ip{w}{F}.
\end{equation}
For sine tests write $b_k=b_{e_k}$ and $b_g=b_{e_g}>0$. The
positivity follows from $e_g>0$ in $\Omega$ and $w\ge0$, $\|w\|_2=1$.
In dimension four $H_1$ is the ordinary Dirichlet Laplacian. For
$0<\alpha<1$, \eqref{eq:form} defines the restricted fractional
Laplacian. We call the principal direction the non-Gaussian mode and the
directions in $w^\perp$ the Gaussian modes.
Write $g=g_n\asymp_{d,R}(\log V)^{-1}$ for the terminal quartic
coefficient at the matched point. Let $W_4$ be a real random variable
with density proportional to $e^{-x^4/4}$. Its Laplace transform
and variance are
\begin{equation}
 \Phi_4(t)=\frac{\int_{\R}e^{tx-x^4/4}\,dx}
                  {\int_{\R}e^{-x^4/4}\,dx},\qquad
 v_4=\E W_4^2=\frac{2\Gamma(3/4)}{\Gamma(1/4)}.
 \label{eq:quartic-transform}
\end{equation}

\begin{theorem}[Susceptibility and exponential moments] \label{thm:tests}
Let $F\in C^\infty(\overline\Omega;\R)$ vanish on the faces and set
$X=S_n(f_n(F))$. With the positive spin amplitudes in
\eqref{eq:amplitudes}, the following asymptotics hold as $n\to\infty$
at fixed sufficiently large $R$. The exponential-moment estimates
are uniform for $\theta$ in bounded real intervals.
\begin{itemize}
\item \emph{Non-Gaussian scaling.} If $b_F\ne0$, then
\begin{equation}
 \chi_n(F)=\frac{n^s}{\sqrt g}\, A_n^{\rm NG}b_F^2(1+o(1)).
 \label{eq:ng-susceptibility}
\end{equation}
\begin{equation}
 \E_{\rm sp}\exp\!\left(\theta g^{1/4}n^{-s/2}X\right)
   =\Phi_4\!\left(\theta b_F\sqrt{A_n^{\rm NG}/v_4}\right)+o(1).
 \label{eq:ng-laplace}
\end{equation}
\item \emph{Gaussian scaling.} If $b_F=0$, then
\begin{equation}
 \chi_n(F)=n^s A_n^{\rm G}\ip{F}{\Ga F}+o(n^s).
 \label{eq:g-susceptibility}
\end{equation}
\begin{equation}
 \E_{\rm sp}\exp\!\left(\theta n^{-s/2}X\right)
   =\exp\!\left(\tfrac12\theta^2 A_n^{\rm G}\ip{F}{\Ga F}\right)+o(1).
 \label{eq:g-laplace}
\end{equation}
\end{itemize}
\end{theorem}

This theorem is an analogue of \cite[Theorem~1.9]{park-torus}.
The moments here are computed at $\tau_{\Lambda_n}^{\rm pc}$.
Appendix~\ref{app:shift} shows that this inverse temperature converges to
the bulk critical point and lies at least $c_{d,R}n^{-d/2}$ above it
for all sufficiently large $n$.

For $F\ne0$, division by the spin standard deviation removes the
possibly $n$-dependent amplitude. Theorem~\ref{thm:tests} then gives
\begin{equation}
 \frac{X}{\sqrt{\chi_n(F)}}\ \xrightarrow{\ \mathrm{law}\ }\
 \begin{cases}
 W_4/\sqrt{v_4}, & b_F\ne0,\\
 N(0,1), & b_F=0 .
 \end{cases}
 \label{eq:standardized-laws}
\end{equation}
The symmetry of $W_4$ removes the sign of $b_F$ in the first limit.


% The Gaussian coefficient is strictly positive for $F\ne0$ in $w^\perp$.
% The corresponding spin sums with bounded second moments are
% $g^{1/4}n^{-s/2}S_n(f_n(F))$ in \eqref{eq:ng-susceptibility} and
% $n^{-s/2}S_n(f_n(F))$ in \eqref{eq:g-susceptibility}. For every fixed sine
% index $k$, take $F=e_k$.

% The spin amplitude cancels in the ratio of two mode susceptibilities,
% giving
% \[
%  \lim_{n\to\infty}\frac{\chi_{k,n}}{\chi_{k_g,n}}
%     =\frac{|\int_\Omega w(z)e_k(z)\,dz|^2}
%            {|\int_\Omega w(z)e_g(z)\,dz|^2}.
% \]
% When the numerator vanishes, \eqref{eq:g-susceptibility} gives the leading term
% at the smaller scale.

Two tests can couple to the same non-Gaussian mode. To determine their covariance
to accuracy $o(n^s)$, we retain the finite-volume susceptibility
$\chi_{k_g,n}$ in the principal term. Its relative asymptotic in
\eqref{eq:ng-susceptibility} alone allows an error larger than $n^s$. Put
\begin{equation}
 j_F=F-\frac{b_F}{I_4}w^3,\qquad
 \mathcal B(F,H)=\ip{j_F}{\Ga j_H},\qquad r_F=b_F/b_g.
\end{equation}
The map $F\mapsto j_F$ is the projection onto $w^\perp$ along $w^3$,
since $\ip{w}{w^3}=I_4$. The direction $w^3$ comes from the first
derivative of the quartic potential.
\begin{theorem}[Mixed moments at the Gaussian scale]\label{thm:mixed}
For fixed smooth $F,H$ vanishing on the faces as in Theorem~\ref{thm:tests},
\begin{equation}
 \Cspin(F,H)=r_F r_H\,\chi_{k_g,n}
 + n^s A_n^{\rm G}
    \left[\mathcal B(F,H)-r_F r_H \mathcal B(e_g,e_g)\right]+o(n^s).
 \label{eq:mixed}
\end{equation}
In particular, if $b_F=b_H=0$, then
$n^{-s}\Cspin(F,H)=A_n^{\rm G}\ip{F}{\Ga H}+o(1)$.
\end{theorem}

\subsection{Related works}
\label{sec:literature}

\paragraph{Temperature shift and rounding.}
Fisher and Barber distinguished the displacement of a finite-system
transition from the width of the temperature interval over which it is
rounded \cite{fisher-barber}. The distinction matters for a free box
when the bulk critical temperature lies outside the finite-volume
critical window. For periodic systems, Br\'ezin and Zinn-Justin
integrated out the nonconstant modes to obtain an effective Hamiltonian
for the spatial average \cite{brezin-zinnjustin}. Its quartic term
controls the fluctuations of the average even when the renormalised
coupling tends to zero. Rudnick, Gaspari and Privman used standing
waves to analyse the shift of the transition and the susceptibility in
a free Ising box above four dimensions \cite{rudnick}. In simulations
of the five-dimensional Ising model, Wittmann and Young found free-box
susceptibility scaling as $L^2$ at the bulk critical temperature and as
$L^{5/2}$ near the finite-volume pseudocritical temperature, where $L$
denotes the side length \cite{wittmann}.

\paragraph{Role of the Fourier modes on the finite-size scaling.}
Wittmann and Young proposed that the all-odd sines inherit the scaling
of the uniform magnetization through their nonzero spatial means.
Their mode susceptibility measurements included the index $(2,1,1,1,1)$,
whose spatial mean vanishes \cite[Sec.~V.B]{wittmann}.
Flores-Sola et al.\ gave the prediction separately for each boundary
condition and temperature \cite[Table~I]{flores}. Write $s$ for their
dispersion exponent $\sigma$, with $s=2$ for short-range interactions.
For $d>2s$ and free boundaries at $T_L$, they assign all-odd indices
to the $Q$ sector, with susceptibility of order $L^{d/2}$. An even
entry places the index in the Gaussian ($G$) sector, with susceptibility
of order $L^s$. At the bulk critical temperature their free-boundary
row gives Gaussian scaling for both classes. Their figures compare
the total magnetization and the lowest $G$ modes in five-dimensional
short-range and one-dimensional long-range Ising models.
The total magnetization alone cannot distinguish the contributions of
two individual all-odd sines.

Honchar, Berche, Holovatch and Kenna subsequently compared the two
all-odd indices $(1,1,1,1,1)$ and $(3,1,1,1,1)$ on a five-dimensional
free box \cite[Fig.~7]{honchar}. They took sine coefficients of a
magnetization profile formed by aligning each spin configuration with
the sign of its total magnetization. The fitted slopes were $-1.26$
and $-1.41$, respectively, and the authors left the asymptotic
interpretation of the mode data open. The quantity $\chi_{k,n}$ studied
here is the second moment of the sine coefficient of a spin
configuration before alignment.

\paragraph{Logarithms at the upper critical dimension.}
At $d=2s$ the two powers printed in Table~I coincide. The quartic
fluctuation can nevertheless be detected by its multiplicative
logarithm. For periodic long-range Ising models, Luijten and Bl\"ote
found the predicted logarithmic correction to the uniform susceptibility
at the marginal decay exponent \cite{luijten}. Lv, Xu, Sun, Chen and
Deng measured the distinction directly in periodic four-dimensional
$O(N)$ models with $N=1,2,3$. The uniform susceptibility grew as
$L^2\sqrt{\log L}$ and the sampled nonzero Fourier modes grew as $L^2$
\cite{lv}. For the weakly coupled $|\varphi|^4$ model on a Euclidean
torus, Park proved the distinction between the zero and nonzero Fourier
susceptibilities, including the marginal logarithm
\cite[Corollary~1.10]{park-torus}. Translation invariance makes the
constant function the principal mode in these periodic examples.
The free-box calculation requires its spatial profile.

\paragraph{Rigorous results for free boxes and critical windows.}
For the nearest-neighbour Ising model in $d>4$, Camia, Jiang and
Newman deduced order $L^2$ for the free-box susceptibility at the bulk
critical temperature from mean-field decay of the infinite-volume
two-point function \cite{camia-jiang-newman}. The bounds of
Duminil-Copin and Panis establish that decay for every $d\ge5$
\cite{duminil-panis}.
For the hierarchical $|\varphi|^4$ model, Michta, Park and Slade
proved that the spatial average has the same quartic scaling law in
the periodic and free-boundary windows about their respective effective
critical points. Above four dimensions the free-boundary shift has
order $L^{-2}$ and the window width has order $L^{-d/2}$
\cite{michta-park-slade}. Park and Slade also obtained the two-point
plateau in these windows, including the logarithmic correction in four
dimensions \cite{park-slade}. For Euclidean lattices, Liu, Park and
Slade connect the periodic susceptibility and the plateau to estimates
on the infinite lattice. Their Conjecture~6 concerns the susceptibility
and two-point profiles in a shifted free-boundary window above the upper
critical dimension \cite{liu-park-slade}. At the matched temperature
studied here, Theorems~\ref{thm:tests} and~\ref{thm:mixed} give the
spatial profile of the enhanced fluctuation and the covariance on its
orthogonal complement.


\subsection{Sine modes and the orthogonality condition}\label{sec:sines}

For every $n$ and every lattice test $f$, Fourier expansion in the
Dirichlet sine basis gives the exact identities
\begin{align}
 M^{(n)}_{kl}
  &=\sum_{x,y\in\Lambda_n}\phi_{k,n}(x)G_{\rm sp}(x,y)\phi_{l,n}(y),
 &\chi_{k,n}&=M^{(n)}_{kk},\label{eq:matrix}\\
 f&=\sum_k d_k\phi_{k,n},\qquad d_k=\ip{\phi_{k,n}}{f},
 &\E_{\rm sp}[S_n(f)^2]&=\sum_{k,l}d_k M^{(n)}_{kl}d_l .
 \label{eq:exact}
\end{align}
If $f\perp\phi_{g,n}$ the sum in \eqref{eq:exact} ranges over
$k,l\ne k_g$. Reflection of coordinate $i$ about $(n+1)/2$ multiplies
the $k$th sine by $(-1)^{k_i+1}$ and preserves the spin law.
Consequently $M^{(n)}_{kl}=0$ whenever $k_i,l_i$ have different parity
for some $i$. The ground state $w$ is invariant under these reflections
by simplicity. Thus $b_k=0$ if any $k_i$ is even. Write
$F_k=\ip{e_k}{F}$ for the continuum sine coefficients and
$|k|^2=\sum_{i=1}^d k_i^2$.

\begin{corollary}[Tests orthogonal to the lowest sine]\label{cor:complement}
In $d=4$, we have $w=e_g$, $I_4=(3/2)^d$, and $b_g=1$. For each fixed
$k\ne k_g$ the susceptibility is
\begin{equation}
 \chi_{k,n}=\frac{n^2 A_n^{\rm G}}
                   {\pi^2(|k|^2-d)}+o(n^2).
\end{equation}
For all fixed $k,l\ne k_g$, including distinct indices in a degenerate
eigenspace,
\begin{equation}
 M^{(n)}_{kl}=n^2 A_n^{\rm G}
       \frac{\delta_{kl}}{\pi^2(|k|^2-d)}+o(n^2).
\end{equation}
If $F\perp e_g$ is as in Theorem~\ref{thm:tests}, its susceptibility is
\begin{equation}
 \chi_n(F)=n^2 A_n^{\rm G}
       \sum_{k\ne k_g}\frac{F_k^2}{\pi^2(|k|^2-d)}+o(n^2).
 \label{eq:d4sum}
\end{equation}

In dimensions one through three take the same class of fixed smooth
tests $F\perp e_g$ vanishing on the faces, and write
$b_F=\sum_{k\ne k_g} b_k F_k$. If $b_F\ne0$, its leading
susceptibility is \eqref{eq:ng-susceptibility}, with the square of this sum.
If $b_F=0$, its leading susceptibility is \eqref{eq:g-susceptibility}, with
\begin{equation}
 \ip{F}{\Ga F}=\lim_{N\to\infty}
 \sum_{\substack{k,l\ne k_g\\ |k|_\infty,|l|_\infty\le N}}
      F_k F_l\ip{e_k}{\Ga e_l}.
 \label{eq:fracsum}
\end{equation}
For each fixed $k$ with $b_k=0$, the Gaussian coefficient in
\eqref{eq:g-susceptibility} is explicitly the spectral sum
$\sum_{j\ge2}\ip{w_j}{e_k}^2/(\mu_j-\mu_1)$, in any real orthonormal
eigenbasis of $H_\alpha$ ordered so that $H_\alpha w_j=\mu_jw_j$
and $w_1=w$. Equation~\eqref{eq:mixed}
with $F=e_k$, $H=e_l$ gives the corresponding off-diagonal entries.
\end{corollary}

\begin{proof}
The eigenpairs of $H_1$ are $(\pi^2|k|^2,e_k)$. Inserting them into
\eqref{eq:g-susceptibility} and \eqref{eq:mixed} gives the four-dimensional
assertions. For fractional powers, Parseval gives the stated expression
for $b_F$ and the spectral theorem gives the sum over $w_j$. Since
$\|\Ga\|=(\mu_2-\mu_1)^{-1}$, rectangular sine truncations converge
in its quadratic form for every $F\in L^2$. The susceptibility
asymptotic for a smooth $F$ is proved directly for its sampled source
in Section~\ref{sec:proof}. The infinite series in \eqref{eq:d4sum}
and \eqref{eq:fracsum} therefore require only convergence in the
bounded continuum form.
\end{proof}

\paragraph{A counterexample to the all-odd criterion at $d=d_c$.}
In four dimensions take $k_*=(3,1,1,1)$. Integrating the sines gives
$\int_\Omega e_{k_*}=\tfrac13\int_\Omega e_g>0$.
Both sines thus couple to a uniform applied field. Their projections
onto the non-Gaussian mode are $b_g=1$ and $b_{k_*}=0$, since $w=e_g$.
The Dirichlet eigenvalue gap is
$\pi^2(|k_*|^2-|k_g|^2)=8\pi^2$. Corollary~\ref{cor:complement} gives
\begin{equation}
 \chi_{k_*,n}=\frac{n^2 A_n^{\rm G}}{8\pi^2}+o(n^2),
 \qquad
 \frac{\chi_{k_*,n}}{\chi_{k_g,n}}\asymp(\log n)^{-1/2}
 \longrightarrow 0 .
 \label{eq:odd-pair}
\end{equation}
Thus two indices in the all-odd column of Table~I have different
susceptibility scales at the upper critical dimension. This disproves
the extension of its free-boundary $Q$ assignment to $d=2s$.
Testing the power-law assignment in the table's stated range requires
a mode estimate for $d>2s$.

For $0<\alpha<1$, reflection symmetry still rules out enhancement
whenever an index has an even entry. The fractional ground state has infinitely many
nonzero higher sine coefficients, as proved in
Appendix~\ref{app:sine-support}. Choose two such indices $k,l$.
The test $F=b_l e_k-b_k e_l$ is orthogonal to both $e_g$ and $w$.
Its susceptibility has order $n^s$, although the two diagonal
contributions in \eqref{eq:exact} each have order $n^s/\sqrt g$.
Their leading terms cancel with the mixed moment.

% \paragraph{Sampling an orthogonal test.}
% For $F\perp e_g$ as above set
% \begin{equation}
%  f_n^\perp=f_n(F)-\phi_{g,n}\ip{\phi_{g,n}}{f_n(F)}.
% \end{equation}
% This gives exact orthogonality. The subtracted coefficient is the
% product trapezoidal sum for $e_gF$, since the integrand vanishes on
% every face. Its integral is zero, so the coefficient is $O_F(n^{-2})$
% by the one-dimensional trapezoidal error in each coordinate. For a
% finite sum of higher sines it is zero once all are represented. For any
% sequence of tests $v_n$ with $\sup_n\sqrt V\,\|v_n\|_\infty<\infty$, entrywise
% nonnegativity of $G_{\rm sp}$ and the uniform-source theorem give
% $\E_{\rm sp}[S_n(v_n)^2]=O(n^s/\sqrt g)$.
% Cauchy--Schwarz therefore bounds the change in each fixed covariance
% caused by this projection by $O_F(n^{-2}n^s/\sqrt g)=o(n^s)$.
% For exponential moments, $|S_n(\phi_{g,n})|\le\sqrt V=n^s$ in every
% configuration. At either normalization the change in the exponent
% is therefore $O_{F,D}(n^{s/2-2})=o(1)$ for $|\theta|\le D$, so
% \eqref{eq:ng-laplace} and \eqref{eq:g-laplace} also hold for
% $S_n(f_n^\perp)$. More
% generally, if $b_F=0$, the $L^2$ norms of the spin sums tested by
% $f_n(F)$ and $\phi_{g,n}$ are respectively $O(n^{s/2})$ and
% $O(n^{s/2}g^{-1/4})$. Hence adding $\delta_n\phi_{g,n}$ to $f_n(F)$
% changes its variance by $o(n^s)$ whenever $|\delta_n|=o(g^{1/4})$.
% For a fixed $F\perp w$, Lemma~\ref{lem:spectrum} gives a power-law
% bound on the error in sampling the principal profile. The finite-volume
% identities \eqref{eq:matrix}--\eqref{eq:exact} apply to every lattice
% test.

\section{Proof of the susceptibility formulas}\label{sec:proof}

The origin of the two scales can be seen from a single field
configuration $z=au$, with $u$ a principal profile of order
$V^{-1/2}$ per site. Its quartic energy has order $ga^4/V$, so the
second moment of $a$ has order $(V/g)^{1/2}=n^s/\sqrt g$. Projecting
away that direction leaves a covariance of order $n^s$. To recover its
coefficient for spins we follow a spatially varying magnetic source
through the Gaussian integrations and retain the derivative of the
quartic potential.

\subsection{The principal profile and the reduced Green operator}
\label{sec:spectrum}

Use the isometric cell embedding \(U_n\) that sends a lattice vector \(v\)
to the function with value \(n^{d/2}v_x\) on
\(\prod_i((x_i-1)/n,x_i/n)\), extended by zero outside \(\Omega\).
Its adjoint takes normalized cell averages. Write
\(\mathcal D(u)=\mathcal E_\alpha(u,u)\) for the form in \eqref{eq:form}.
The first estimate below controls the error in the overlap with $w$.
A rate is needed because its squared overlap is multiplied by
$n^s/\sqrt g$. The second estimate identifies the coefficient on
the $n^s$ scale.

\begin{lemma}[Principal profile and reduced covariance]\label{lem:spectrum}
Take the prescribed Gaussian reference of \cite[Section~2]{base}, with
\(B-A=\zeta J_n\), bulk row sum \(p=\sum_z A_\infty(z)>0\), and
\(G=a_R(A^{-1}-c'B^{-1})^{-1}\). Let \(\ell\) be the positive
solution of \(A\ell=\eta_1 B\ell\) for the largest generalized
eigenvalue \(\eta_1=\eta_{1,n}\), normalized by
\(\ell^{\mathsf T}B\ell=1\). Put
\begin{equation}
 u_*=B\ell,\qquad
 G_h=G-(\ell^{\mathsf T}G\ell)u_*u_*^{\mathsf T}.
\end{equation}
At fixed \(R\), for \(\delta=s/16\),
\begin{equation}
 \|U_n u_*-\sqrt p\,w\|_{L^2}=O_R(n^{-\delta}).
\end{equation}
For any lattice sequences \(v_n,z_n\) with \(U_n v_n\to v\) and
\(U_n z_n\to z\) in \(L^2(\Omega)\),
\begin{equation}
 n^{-s}\langle v_n,G_h z_n\rangle
 \longrightarrow \frac{a_Rp^2}{\zeta\kappa_R}
                      \langle v,\mathcal R_\alpha z\rangle .
\end{equation}
\end{lemma}

Proposition~2.7 of \cite{base} controls the first generalized
eigenvalue and the gap in the box. Appendix~\ref{app:spectral}
compares the walk form with \eqref{eq:form} to obtain the rate for
the principal eigenvector and the limit of the reduced covariance.
For smooth $F$ the sampling satisfies
$\|U_nf_n(F)-F\|_2=O_F(n^{-1})$. Indeed the sample point and any
point in its cell differ by $O(n^{-1})$, and the ratio of the two
normalizing factors is $(n/(n+1))^{d/2}=1+O(n^{-1})$.


\subsection{Smooth sources and the spin normalization}\label{sec:sources}

The microscopic split replaces a spin source $f$ by a convolution of
$f$ restricted to the cube. For a signed source it is important to
bound the boundary error on the $n^s$ scale, even if the variance of
the source itself is larger. Lemma~\ref{lem:sources} gives that bound.

Let $\psi$ denote the full canonical field in the coordinates used at
the final integration. Its Gaussian reference has covariance $\Sigma$.
In the original coordinates the reference covariance is $G$ of
Lemma~\ref{lem:spectrum}. The change of field is given by
\begin{equation}
 \Sigma=tG+\varepsilon a_R A,\qquad S=G\Sigma^{-1},\qquad
 h_{\rm phys}=a_R^{-1/2}S\psi .
\end{equation}
For $d\le3$ this gives $\Sigma=G$ and $S=I$. For $d=4$ it is the
change of variables of \cite[Lemma~7.10]{base} at the prescribed match.
Write $\E_{\rm full}$ for expectation under the interacting law of
$\psi$. For a covariance $D$, write $\E_D$ for expectation under
the centred Gaussian of covariance $D$.

In the reference measure write $\psi=\psi_<+\psi_T$ as a sum of
independent Gaussian fields of covariances $C_<$ and $C_T$.
Thus $\Sigma=C_<+C_T$. The renormalisation group integrates $\psi_<$ first. Regressing
$\psi_T$ on its principal coordinate $a=\ell^{\mathsf T}\psi_T$ gives
\begin{align}
 \sigma_T^2&=\ell^{\mathsf T}C_T\ell,\qquad
 u=C_T\ell/\sigma_T^2,\qquad \psi_T=au+Y, \\
 C_h&=C_T-\sigma_T^2uu^{\mathsf T},\qquad
 K=C_<+C_h=\Sigma-\sigma_T^2uu^{\mathsf T}. \label{eq:cov-split}
\end{align}
Here $a$ and $Y$ are independent under the reference measure, with
laws $N(0,\sigma_T^2)$ and $N(0,C_h)$. Thus $u$ is the terminal
regression profile and $u_*=B\ell$ is the profile for the full
reference. In $d=4$, $a$ in \eqref{eq:cov-split} is the coordinate
of the transformed field at cut $k$. The estimates below use its
distribution at the temperature selected by $a_0$. The terminal
test norm $E_n$ in \cite[Section~7.1]{base} is the weighted
difference norm specified in Appendix~\ref{app:slices}.

\begin{lemma}[Source directions and the return to spins]\label{lem:sources}
Let \(F,H\) be smooth on \(\overline\Omega\) and vanish on
\(\partial\Omega\). Put
\(f=f_n(F)=(n+1)^{-d/2} F(\cdot/(n+1))\) and \(h=f_n(H)\).
At fixed \(d,R\) the source and covariance satisfy
\begin{equation}
 \|C_<\|_{2\to2}+\|K\|_{2\to2}=O_R(n^s),\qquad
 \|C_< f\|_{E_n}+\|Kf\|_{E_n}=O_{R,F}(1).
\end{equation}
For the same tests, the covariance of the spins satisfies
\begin{equation}
 \mathcal C_n(F,H)=\frac{1}{\beta a_R m^2}
 \E_{\rm full}\!\left[\ip{f}{\psi}\ip{\psi}{h}\right]+o_{R,F,H}(n^s).
\end{equation}
Set $c_{\rm sp}=1/(m\sqrt{\beta a_R})$. For every
fixed $D<\infty$, uniformly for real $|\theta|\le D n^{-s/2}$,
\begin{equation}
 \E_{\rm sp}e^{\theta S_n(f)}
   =(1+o(1))\,\E_{\rm full}e^{\theta c_{\rm sp} f^{\mathsf T}\psi}.
 \label{eq:spin-laplace-transfer}
\end{equation}
\end{lemma}

\begin{proof}
Appendix~\ref{app:slices} proves the bounds on $C_<$ and $K$ by
summing the covariance slices with a smooth source. To transfer the
covariance to spins, use the exact Hubbard--Stratonovich representation
in physical units,
\begin{equation}
 h_{\rm phys}=a_R^{-1/2}S\psi
 =\sqrt\beta A\sigma+\eta_{\rm phys},\qquad
 \eta_{\rm phys}\sim N(0,A)
\end{equation}
in law, with independent noise. If \(v=S^{-\mathsf T}f\) and
\(v_h=S^{-\mathsf T}h\), its covariance identity is
\begin{equation}
 a_R^{-1}\E_{\rm full}[\ip{f}{\psi}\ip{\psi}{h}]
   =\beta(Av)^{\mathsf T}G_{\rm sp}(Av_h)+v^{\mathsf T}Av_h.
\end{equation}
The noise term \(v^{\mathsf T}Av_h\) is \(O_{R,F,H}(1)\), since \(A\) and
\(S^{-\mathsf T}\) are uniformly bounded on \(\ell^2\) at fixed \(R\).

For some \(0<\gamma<\min\{1,s\}\) the clipped rows satisfy
\(\|Af-pf\|_\infty=O_{R,F}(n^{-s-\gamma})\).
Indeed the zero extension of \(F\) is Lipschitz, and the bulk
kernel of \(A\) has a finite absolute \(\gamma\)-moment. For
the fractional split the tail expansion has coefficients
\(O(j^{-1-s/2})\), and the \(j\)-step spread-out walk
has \(\gamma\)-moment at most \(C_R j^{\gamma/2}\). In dimension
four the kernel has finite range. Transposing
$S^{-1}=\Sigma G^{-1}=I-\varepsilon c'AB^{-1}$ gives
\(S^{-\mathsf T}=I-\varepsilon c' B^{-1}A\). Since \(B\) has bulk
row sum \(p\), finite range, and uniformly bounded inverse in
absolute row norm, \(Bf=pf+O_{R,F}(n^{-s-1})\) pointwise implies
\(B^{-1}f=p^{-1}f+O_{R,F}(n^{-s-1})\). The preceding bounds give
\begin{equation}
 Av=m f+O_{R,F}(n^{-s-\gamma})\quad\text{in supremum norm}.
\end{equation}
The same identity in lower dimensions uses \(m=p\).

Ferromagnetic positivity and the uniform susceptibility bound in
\cite[Lemma~7.5 and Theorem~1.1]{base} give \(G_{\rm sp}(x,y)\ge0\)
and
\(\sum_{x,y}(G_{\rm sp})_{xy}=O_R(V n^s\sqrt{\log n})\).
Consequently each error term in replacing \((Av)^{\mathsf T}G_{\rm sp}(Av_h)\)
by \(m^2 f^{\mathsf T}G_{\rm sp}h\) is
\(O_{R,F,H}(n^{s-\gamma}\sqrt{\log n})=o(n^s)\).
The factors \(\beta\) and \(m\) are bounded above and away from
zero at fixed \(R\). This proves the covariance formula.

For exponential moments we need a bound on the source error for
each spin configuration. Let $r=Av-mf$. In $d\le3$ choose
$s/2<\gamma<\min\{1,s\}$ in the clipped-row estimate above. Since
$V=n^{2s}$, we have $\|r\|_1=O(n^{s-\gamma})=o(n^{s/2})$.
In $d=4$ the kernels $A$ and $B$ are symmetric and have finite
range at fixed $R$. In an interior row, Taylor expansion of $F$ and
cancellation of the linear term give errors $O_{R,F}(n^{-s-2})$
in $Af-pf$ and $Bf-pf$. In the fixed-width boundary collar,
$F|_{\partial\Omega}=0$ gives errors $O_{R,F}(n^{-s-1})$ on
$O_R(n^{d-1})$ sites. Both $\ell^1$ errors are therefore
$O_{R,F}(n^{s-2})=O_{R,F}(1)$. The identity
$AB^{-1}Af-pf=(A-pI)f+AB^{-1}(A-B)f$ gives the same bound for
$r$ because $A$ and $B^{-1}$ have bounded absolute column sums.
Thus $|\theta r^{\mathsf T}\sigma|\le|\theta|\|r\|_1=o(1)$ at
the stated scales in every dimension.

The independent Gaussian noise in the Hubbard--Stratonovich identity
gives the exact relation
\begin{equation}
 \E_{\rm full} e^{\theta c_{\rm sp} f^{\mathsf T}\psi}
 =\exp\!\left(\frac{\theta^2 v^{\mathsf T}Av}{2\beta m^2}\right)
   \E_{\rm sp}\exp\!\left(\frac{\theta}{m}(Av)^{\mathsf T}\sigma\right).
\end{equation}
The first factor tends uniformly to one since $v^{\mathsf T}Av=O(1)$.
The pointwise error estimate gives \eqref{eq:spin-laplace-transfer}.
\end{proof}

\subsection{Quartic moments of the principal coordinate}\label{sec:terminal}

Let $F_{\rm in}$ be the density of the full interacting field relative
to its Gaussian reference, up to a constant. Integration of the
earlier scales gives $F_T(z)=\E_{C_<}[F_{\rm in}(z+\psi_<)]$, up
to a positive field-independent factor. Thus the interacting law of
$\psi_T$ is the Gaussian law of covariance $C_T$ weighted by
$F_T(\psi_T)/Z$, where $Z=\E_{C_T}[F_T(\psi_T)]$.
Write $D_v$ for the directional derivative in the real field
direction $v$. The principal coordinate has quartic scale
$a_*=(V/g)^{1/4}$, corresponding to the site-field scale
$h_*=a_*/\sqrt V$. With the normalized filter $T$ described after
Theorem~\ref{thm:mixed}, the terminal estimates give
\begin{gather}
 \sigma_T^2\asymp_R n^s/g,\qquad
 \sqrt V\,\|u\|_{E_n}+\|\ell\|_2\le C_R,\qquad
 \E_{C_h}\exp(c_R n^s\|Y\|_{E_n}^2)\le C_R,\label{eq:terminaldata}\\
 F_T(z)=\exp[-V_4(z)]+\mathcal E(z),\qquad
 V_4(z)=\frac{g}{4}\sum_x(Tz)_x^4 .
\end{gather}
The weight $F_T$ is real and nonnegative. For $j=0,1,2$ its remainder
satisfies the Gaussian envelope
\begin{equation}
 \E_{C_h}\left|D_{v_1}\cdots D_{v_j}\mathcal E(a_* y u+Y)\right|
 \le C_R\epsilon_n e^{-c_R y^2}
              \prod_{i=1}^j\frac{\|v_i\|_{E_n}}{h_*},
 \qquad \epsilon_n\to0.
 \label{eq:envelope}
\end{equation}
Appendix~\ref{app:terminal-inputs} specifies the matched-trajectory
estimates that give \eqref{eq:terminaldata}--\eqref{eq:envelope} in
each dimension. The same estimate for $e^{-V_4}$ holds with $\epsilon_n$ replaced by
one. In the proof of Proposition~7.4 of \cite{base}, the $L^2$
regulator gives the Gaussian decay in $y$ because
$\ell^{\mathsf T}(a_*yu+Y)=a_*y$. The remaining regulator is
integrable by \eqref{eq:terminaldata}. For the pure quartic, the
bounded $\ell^4$ inverse of $T$ gives $V_4(z)\ge c g\|z\|_2^4/V$.
Its field derivatives give polynomials which are absorbed by a
fraction of the quartic. These envelopes also permit any fixed
polynomial weight in $y$ under integration.

The coefficient of the enhanced susceptibility comes from the second
moment of the scalar quartic density. Its fourth moment appears when
the source differentiates the quartic and determines the correction
from $F$ to $j_F$ in the mixed covariance.

\begin{lemma}[Scalar moments and the differentiated quartic]\label{lem:quartic}
Write $m_{2,n}=Z^{-1}\E_{C_T}[a^2F_T(\psi_T)]$ for the second moment of the interacting principal coordinate. Define
\begin{equation}
 q_n=V\sum_x(Tu)_x^4,\qquad
 d_n=\frac{V}{q_n}T^{\mathsf T}[(Tu)^3],
\end{equation}
where the cube in brackets is pointwise. Then
\begin{align}
 q_n&\longrightarrow p^2 I_4,&
 m_{2,n}&=\frac{n^s}{\sqrt g}\,
       \frac{2\Gamma(3/4)}{\Gamma(1/4)\sqrt{q_n}}(1+o(1)),
 \label{eq:gamma}\\
 \ip{u}{d_n}&=1,& U_n d_n&\longrightarrow
              \frac{w^3}{\sqrt p\,I_4}\quad\text{in }L^2 .
 \label{eq:cubic}
\end{align}
For real directions $v$ with $\|v\|_{E_n}=O(1)$,
\begin{equation}
 Z^{-1}\E_{C_T}[aD_v F_T]=-\ip{d_n}{v}+o(n^s),\qquad
 Z^{-1}\E_{C_T}|D_v^2 F_T|=O_R(\sqrt g\,n^s).
 \label{eq:direction}
\end{equation}
The errors are uniform on bounded sets in the indicated norm.
\end{lemma}

\begin{proof}
The full covariance in the natural field is $\Sigma=tG+\varepsilon a_RA$.
Because $A\ell=\eta_1 B\ell$ and $G\ell=(\ell^{\mathsf T}G\ell)u_*$,
its full regression profile is exactly $u_*=B\ell$. Subtracting the
preterminal column and its variance gives
\begin{equation}
 u-u_*=
   \frac{(\ell^{\mathsf T}C_<\ell)u_*-C_<\ell}{\sigma_T^2},
 \qquad \|u-u_*\|_2=O_R(g)
\end{equation}
by Lemma~\ref{lem:sources} and \eqref{eq:terminaldata}. Hence
$U_nu\to\sqrt p\,w$ in $L^2$ by Lemma~\ref{lem:spectrum}.
The uniform bound on $U_nu$ also gives $w\in L^\infty$.
At fixed $R$ the normalized filter has bounded absolute row and column
sums and microscopic range, with row sum one away from its boundary
collar. The bulk filter is symmetric, and the same properties hold
for its transpose. In $d\le3$ the filter is the normalized finite
Neumann sum in the proof of \cite[Proposition~7.2]{base}. In $d=4$
it is the normalized short-walk filter of \cite[Lemma~7.10]{base}.
On samples of a smooth compactly supported function
these filters converge strongly to the identity by translation
continuity or the first moment bound. Density and uniform operator
bounds give the same assertion on any strongly convergent embedded
sequence, including along the allowed counterterms. Thus
$U_nTu\to\sqrt p\,w$ strongly in $L^2$ and is uniformly bounded
in $L^\infty$. Taking the fourth power in $L^1$ and the cube in
$L^2$, then applying $T^{\mathsf T}$ to the latter, gives the limits
of $q_n$ and $U_nd_n$. The identity $u^{\mathsf T}d_n=1$ follows
by contraction.

Write $a=a_*y$. On bounded $y$-sets,
$g^{1/4}\|T(a_*yu)\|_4=q_n^{1/4}|y|$, whereas
$g^{1/4}\|TY\|_4=O_{L^j}(g^{1/4})$ for every fixed $j$
by \eqref{eq:terminaldata} and $\|TY\|_4\le C V^{1/4}\|Y\|_\infty$.
Consequently $V_4(a_*yu+Y)-q_n y^4/4\to0$ in $L^1$, uniformly
on these sets. The scalar Gaussian factor
$\exp[-a_*^2y^2/(2\sigma_T^2)]$ tends to one there.
Equation~\eqref{eq:envelope} controls the tails in the integrals
with weights $1,y^2,y^4$. In particular $Z\asymp_R a_*/\sigma_T$.
After dividing by $Z$, the second and fourth moments of $y$ equal
those under the density proportional to $\exp(-q_n y^4/4)$, up to
$o(1)$. The same comparison holds with $e^{-V_4}$ in the numerator.
The second moment of this scalar quartic law is
$2\Gamma(3/4)/(\Gamma(1/4)\sqrt{q_n})$, and its fourth moment
is $1/q_n$. This proves \eqref{eq:gamma} and justifies moment
convergence at the normalization used in Theorem~\ref{thm:tests}.

For the first derivative of the pure quartic weight the insertion is
$-g a\sum_x(T(au+Y))_x^3(Tv)_x$. Replacing the cube by
$a^3(Tu)_x^3$ gives a normalized error $o(n^s)$. In detail, a term
with $j\ge1$ factors of $TY$ costs at most
\begin{equation}
 C g|a|^{4-j}V\,\|Tu\|_\infty^{3-j}
                    \|TY\|_\infty^j\|Tv\|_\infty .
\end{equation}
At $a=a_* y$ its Gaussian expectation, with the pure quartic
weight, is bounded by $C_R n^s g^{j/4}(1+|y|^4)e^{-c_R y^4}$.
Here $V_4(au+Y)\ge c_R y^4$ follows by $\ell^{\mathsf T}Y=0$ and
the $\ell^4$ inverse bound. The factor $a_*/\sigma_T$ from
integration cancels against $Z$. Since $ga_*^4=V$, the retained term is
\begin{equation}
 -\frac{V}{q_n}\sum_x(Tu)_x^3(Tv)_x\,
       \frac{q_n}{Z}\E_{C_T}\!\left[(a/a_*)^4e^{-V_4(\psi_T)}\right]
 =-\ip{d_n}{v}+o(n^s).
\end{equation}
Here the last factor tends to one by the quartic fourth moment, and
$|\ip{d_n}{v}|=O_R(\sqrt V)$. The derivative of the remainder
costs $O_R(\epsilon_n a_*/h_*)=o(n^s)$ by
\eqref{eq:envelope}. Two derivatives cost
$O_R(h_*^{-2})=O_R(\sqrt g\,n^s)$ by the same envelope for
$F_T$. This proves \eqref{eq:direction}.
\end{proof}

\subsection{Reconstruction of the spin covariance}

\begin{proof}[Proof of the covariance formulas in Theorems~\ref{thm:tests} and \ref{thm:mixed}]
Set $f=f_n(F)$ and $b=u^{\mathsf T}f$. Exact Gaussian
translation, first in the integrated field and then in the independent
terminal variable $Y$, gives for $\eta\in\R$
\begin{equation}
 \E_{\rm full} e^{\eta f^{\mathsf T}\psi}=
  \frac{e^{\eta^2 f^{\mathsf T}Kf/2}}{Z}
       \E_{C_T}\!\left[e^{\eta ab}F_T(au+Y+\eta Kf)\right].
 \label{eq:translate}
\end{equation}
The integrated field translates the interaction by $\eta C_<f$ and the
terminal residual translates it by $\eta C_h f$. Gaussian translation
holds on the support of $Y$ even when $C_h$ is singular. For fixed $n$
and $\eta$ near zero, a fraction of the Gaussian decay in the terminal
bounds absorbs the translation and the linear source. This gives an
integrable bound for the first two derivatives at zero.

Lemma~\ref{lem:sources} gives $\|Kf\|_{E_n}=O(1)$, so
\eqref{eq:direction} applies with $v=Kf$. Two derivatives of
\eqref{eq:translate} produce $f^{\mathsf T}Kf$, $m_{2,n} b^2$,
$2b Z^{-1}\E_{C_T}[aD_{Kf}F_T]$, and $Z^{-1}\E_{C_T}[D_{Kf}^2 F_T]$.
Completing the square in $K$ yields
\begin{equation}
 \E_{\rm full}(f^{\mathsf T}\psi)^2
 =L_n b^2+(f-bd_n)^{\mathsf T}K(f-bd_n)+o(n^s),
 \qquad L_n=m_{2,n}-d_n^{\mathsf T}K d_n .
 \label{eq:square}
\end{equation}
In fact the cross term from \eqref{eq:translate} is
$2bZ^{-1}\E_{C_T}[aD_{Kf}F_T]=-2b\,d_n^{\mathsf T}Kf+o(n^s)$.
The bounded $\ell^2$ norm of $d_n$ gives $L_n=m_{2,n}+O_R(n^s)$.
Applying \eqref{eq:square} to $f+h$ and $f-h$ and taking one quarter
of the difference gives the mixed moment for $f=f_n(F)$ and $h=f_n(H)$.

Put $f'=f-bd_n$, so that $u^{\mathsf T}f'=0$ exactly. We can
evaluate $K$ on this test without differentiating in the direction
$Kd_n$. Indeed $K=\Sigma-\sigma_T^2 uu^{\mathsf T}$, and
$u_*^{\mathsf T} f'=O_{R,F}(g)$ by the profile comparison in
Lemma~\ref{lem:quartic}. The full Gaussian variance in the principal
direction is $O_R(n^s/g)$. Since $A\ell=\eta_1u_*$, removing that
summand from $\Sigma$ leaves exactly
\begin{equation}
 tG_h+\varepsilon a_R(A-\eta_1 u_*u_*^{\mathsf T}).
\end{equation}
The second term is bounded in operator norm. The omitted full
principal term tested on $f'$ costs $O_R(n^s g)$. The same estimates
apply to bilinear tests. Equations~\eqref{eq:cubic} and
Lemma~\ref{lem:spectrum} now give, with $h'=f_n(H)-
(u^{\mathsf T}f_n(H))d_n$,
\begin{equation}
 U_n f'\to j_F,\qquad
 n^{-s}(f')^{\mathsf T}K h'
 =\frac{t a_Rp^2}{\zeta\kappa_R}\,\mathcal B(F,H)+o(1).
 \label{eq:reduced-limit}
\end{equation}
The $o(1)$ holds for the actual bounded sequence $t=t_n$.

For the scalar term in \eqref{eq:square}, smooth sampling and the
ground-state estimate give
\begin{equation}
 u^{\mathsf T}f_n(F)=\sqrt p\, b_F+
                 O_{R,F}(g+n^{-s/16}).
\end{equation}
Multiplying this error by $L_n=O_R(n^s/\sqrt g)$ still costs
$o(n^s)$. Combining \eqref{eq:square}, \eqref{eq:reduced-limit}, and
Lemma~\ref{lem:sources} therefore gives the single bilinear formula
\begin{equation}
 \Cspin(F,H)=\frac{p L_n}{\beta a_R m^2}\,
        b_F b_H+n^s A_n^{\rm G}\mathcal B(F,H)+o(n^s).
 \label{eq:master}
\end{equation}
Using \eqref{eq:gamma} and $q_n\to p^2 I_4$ gives \eqref{eq:ng-susceptibility}.
When $b_F=0$ the remaining form has $j_F=F$ and gives
\eqref{eq:g-susceptibility}. Taking $F=H=e_g$ in \eqref{eq:master} and eliminating
its scalar coefficient gives \eqref{eq:mixed}.
\end{proof}

\subsection{Exponential moments of the spin field}

\begin{proof}[Proof of the exponential-moment assertions of Theorem~\ref{thm:tests}]
Set $f=f_n(F)$ and $b=u^{\mathsf T}f$. Fix $D>0$ and take
$\theta\in[-D,D]$. Write $\rho_n=g^{1/4}n^{-s/2}$ if $b_F\ne0$
and $\rho_n=n^{-s/2}$ if $b_F=0$. To compute
$\E_{\rm sp}e^{\theta\rho_n X}$, use
\eqref{eq:spin-laplace-transfer} and \eqref{eq:translate}.
The two factors are the Gaussian exponential
$\exp(\eta_n^2 f^{\mathsf T}Kf/2)$ and the translated terminal
integral, with
\begin{equation}
 \eta_n=\theta c_{\rm sp}\rho_n,\qquad v=\eta_n Kf,\qquad
 \xi_n=\eta_n a_* b .
\end{equation}
Lemma~\ref{lem:sources} gives $\|v\|_{E_n}/h_*=O(g^{1/4})$ at both
scales. The profile comparison gives
$b=\sqrt p\,b_F+O(g+n^{-s/16})$. Thus $\xi_n$ stays bounded if
$b_F\ne0$. If $b_F=0$, then
$\xi_n=O_D(g^{-1/4}(g+n^{-s/16}))=o(1)$.

We first evaluate the terminal integral with $a=a_*y$. Set
$\delta_n=\|v\|_{E_n}/h_*$. For $\phi=a_*yu+Y+v$, the identities
$\ell^{\mathsf T}Y=0$ and $\ell^{\mathsf T}u=1$ imply
$\|\phi\|_2/a_*\ge c(|y|-C\delta_n)_+$. The pointwise remainder
bound used in Appendix~\ref{app:terminal-inputs} therefore gives
\begin{equation}
 \E_{C_h}|\mathcal E(a_*yu+Y+v)|\le C\epsilon_n e^{-c y^2},
 \qquad
 e^{-V_4(a_*yu+Y+v)}\le \exp[-c(|y|-C\delta_n)_+^4].
 \label{eq:translated-envelope}
\end{equation}
Indeed the positive exponent in the remainder regulator is bounded by
$C c_h(y^2+h_*^{-2}\|Y\|_{E_n}^2+\delta_n^2)$, where $c_h$ is the
regulator constant of \cite[Assumption~7.1]{base}. Its $y^2$ term is
absorbed by the negative quadratic exponent, as in the proof of
\eqref{eq:envelope}, and \eqref{eq:terminaldata} controls the
expectation of the remaining factor. The quartic bound follows from
$V_4(\phi)\ge c g\|\phi\|_2^4/V$.

On bounded $y$-intervals, $g^{1/4}\|Tv\|_4\le C\delta_n$ and
$g^{1/4}\|TY\|_4=O_{L^4}(g^{1/4})$. Hence
$V_4(a_*yu+Y+v)=q_n y^4/4+o_{L^1}(1)$ there, uniformly on each
such interval. The scalar Gaussian factor tends to one there as well.
Since $|\xi_n|$ is bounded, the bounds in
\eqref{eq:translated-envelope} remain integrable after multiplication
by $e^{\xi_n y}$. Dividing by the unshifted integral for $Z$ cancels
the Gaussian normalization $a_*/(\sqrt{2\pi}\sigma_T)$ and gives
\begin{equation}
 \frac{\E_{C_T}\!\left[e^{\eta_n a b}F_T(au+Y+v)\right]}{Z}
 =\frac{\int_{\R}e^{\xi_n y-q_n y^4/4}\,dy}
        {\int_{\R}e^{-q_n y^4/4}\,dy}+o(1)
 =\Phi_4(\xi_n q_n^{-1/4})+o(1).
 \label{eq:terminal-laplace}
\end{equation}
The denominators are bounded away from zero after cancellation since
$q_n\to p^2 I_4>0$. All tail and compact-interval estimates are
uniform for $|\theta|\le D$.

If $b_F\ne0$, then $\eta_n^2 f^{\mathsf T}Kf=O(\sqrt g)$.
Also $\xi_n=\theta c_{\rm sp} b$, so that
$\xi_n q_n^{-1/4}=\theta b_F\sqrt{A_n^{\rm NG}/v_4}+o(1)$ by
\eqref{eq:amplitudes}. This proves \eqref{eq:ng-laplace}.
If $b_F=0$, the terminal factor in \eqref{eq:terminal-laplace}
tends to one. Since $b=o(1)$ and $\|d_n\|_2=O(1)$, replacing $f$
by $f-bd_n$ in $n^{-s}f^{\mathsf T}Kf$ costs $o(1)$ by
Lemma~\ref{lem:sources}. Now $j_F=F$, and \eqref{eq:reduced-limit}
gives $\eta_n^2 f^{\mathsf T}Kf
=\theta^2 A_n^{\rm G}\ip{F}{\Ga F}+o(1)$. This proves \eqref{eq:g-laplace}.
Finally, for $F\ne0$ the susceptibility formulas give
$\rho_n^2\chi_n(F)\asymp_{R,F}1$, since $\Ga$ is strictly positive
on $w^\perp$ and both amplitudes are bounded above and away from zero.
For fixed $t\in\R$ take $\theta=t/(\rho_n\sqrt{\chi_n(F)})$.
This choice stays in a bounded interval, and the two formulas give
the Laplace transforms of the variance-one laws in
\eqref{eq:standardized-laws}.
\end{proof}

\section{Spin amplitudes}
In the microscopic split of \cite[Section~2]{base}, $A_\infty$ is the
positive convolution kernel used for the Gaussian field and $Q_\infty$
is the remaining finite-range spin interaction. Write $A,Q$ for their
restrictions to $\Lambda_n$, and put $B=I-Q$, retaining the diagonal
of $Q$ as in the Gaussian reference. The temperature and source
normalizations are
\begin{equation}
 B-A=\zeta J_n,\qquad \beta=\tau/\zeta,\qquad
 p=\sum_z A_\infty(z)=1-\sum_z Q_\infty(z)>0 .
\end{equation}
Here $\zeta=(1+e^{-R})^{-1}$ for $d\le3$, and $\zeta=(1+t_R)^{-1}$
for the local split with diagonal correction $t_R$ in
\cite[Lemma~2.1]{base}. Denote the microscopic field exponent $a_*$
of \cite[Section~2]{base} by $\gamma_{\rm sc}\in(1/12,1/10)$.
Then $\lambda=R^{-2d\gamma_{\rm sc}}$
and $a_R=\lambda R^s$. The physical Hubbard--Stratonovich field in
the original split equals the canonical field divided by $\sqrt{a_R}$.

Following Section~7.5 of \cite{base}, we write $c'=c'_n$ for the
original reference parameter called $c$ in its Section~2. It lies below
the pole $c_{*,n}=\lambda_{\max}(B^{-1/2}AB^{-1/2})^{-1}$. In $d\le3$ set
$\varepsilon=0$. In $d=4$, let $\varepsilon=\varepsilon_n$ be the
parameter of the change of field selected together with $\tau$ in
\cite[Corollary~7.13]{base}. Put $t=1-\varepsilon$ and
$m=p(1-\varepsilon c')$. The coefficient $g$ is normalized so that
the terminal quartic is $(g/4)\sum_x(Tz)_x^4$, where the linear
filter $T$ has bulk row sum one. In $d\le3$, $g=g_N$ is the terminal
coupling of Proposition~7.14 of \cite{base}. In $d=4$ the prescribed
coordinate $a_0$ at cut $J$ fixes the temperature through its
marginal density \cite[Eq.~(7.43)]{base}. The proof of Corollary~7.13
then evaluates moments in the transformed field at the earlier cut
$k=J-1$. We use $g=g_k$ there. The
parameters $t,m,\beta$ are bounded above and away from zero at fixed
$R$. The amplitudes are
\begin{align}
 \sigma_R^2&=\sum_z z_1^2D_R(z),&
 \kappa_R&=(\sigma_R^2/2)^\alpha,\\
 A_n^{\rm NG}&=\frac{2\Gamma(3/4)}
 {\Gamma(1/4)\,\sqrt{I_4}\,\beta a_R m^2},&
 A_n^{\rm G}&=\frac{t}{\tau\,\kappa_R(1-\varepsilon c')^2}.
 \label{eq:amplitudes}
\end{align}
These are the amplitudes for the original $\pm1$ spins. For example,
in $d\le3$ the Gaussian amplitude is $1/(\tau\kappa_R)$. All statements
allow the matched parameters to depend on $n$. A fixed numerical
coefficient multiplying $n^s\sqrt{\log V}$ would require convergence
of $A_n^{\rm NG}/\sqrt{g\log V}$, and a fixed Gaussian amplitude
would require convergence of $A_n^{\rm G}$.

\appendix
\section{Spectral estimates on the cube}\label{app:spectral}
\subsection{Convergence of the principal profile and the reduced resolvent}

\begin{proof}[Proof of Lemma~\ref{lem:spectrum}]
The domain of \eqref{eq:form} embeds compactly in \(L^2(\Omega)\) by
the translation criterion. For \(s<2\) the jump form has a
strictly positive kernel between distinct points. Hence each real
minimizer has one sign, since taking its absolute value strictly
decreases the energy if both signs occur on sets of positive measure.
Two nonzero orthogonal nonnegative minimizers would have disjoint
supports and their difference would be a sign-changing minimizer.
This proves simplicity. For \(s=2\) simplicity follows from the
Dirichlet sine expansion.

Put \(b=1/n\) and
\(D_n(v)=b^{-s}\langle v,J_n v\rangle/\kappa_R\).
In rescaled Fourier variables its zero-extension multiplier is
\(j_b(\xi)=b^{-s}(1-\widehat D_R(b\xi))^\alpha/\kappa_R\),
where the walk multiplier is its characteristic function. On the
Brillouin zone it is comparable at fixed \(R\) to \(|\xi|^s\)
since the support is aperiodic and the origin is a nondegenerate quadratic
minimum of \(1-\widehat D_R\). For all \(\xi\),
\(j_b(\xi)\le |\xi|^s\) by \(1-\cos t\le t^2/2\) and cube symmetry.
On \(|\xi|\le b^{-1/2}\), Taylor expansion gives
\(j_b(\xi)=|\xi|^s(1+O_R(b))\).

For \(\|v\|_2+D_n(v)\le C\), lattice translations of its zero
extension obey
\(\|v(\cdot+m)-v\|_2\le C_R|bm|^{s/2}\). This follows from
\(|e^{it}-1|^2\le C|t|^s\) and the multiplier lower bound.
For a translation of the embedded step function the squared norm is
a convex combination of the squared norms for the adjacent lattice
translations, by subdividing each cell. Consequently
\begin{equation}
 \|U_n v(\cdot+y)-U_n v\|_2\le C_R(|y|+b)^{s/2},\qquad
 \|(B-pI)v\|_2=O_R(b^{s/2}).
\end{equation}
The second bound uses the finite-range convolution representing
\(Q=I-B\), with bulk row sum \(1-p\), before restriction.

We give the spectral comparison including multiplicities. If \(z\)
belongs to the continuum form domain, the difference-energy
representation of \(J_n\) gives
\begin{equation}
 D_n(U_n^*z)\le\int_{\mathbb R^d}j_b(\xi)|\widehat z(\xi)|^2\,d\xi
 \le \mathcal D(z).
\end{equation}
Indeed on the full cell grid averaging commutes with each lattice
translation and is contractive, and all jump weights are nonnegative.
The cell averages outside the cube are zero. Also
\(\|U_nU_n^*z-z\|_2\le C b^{s/2}\mathcal D(z)^{1/2}\), by averaging
the squared translation bound over pairs of points in each cell.
Conversely, any unit sequence of bounded \(D_n\) has an embedded
strongly convergent subsequence by the translation estimate and
bounded support. Its limit \(z\) satisfies
\(\mathcal D(z)\le\liminf D_n(v)\). Indeed on each fixed Fourier
ball the transform of \(U_n v\) equals the normalized lattice
transform times the cell multiplier, a phase times a product of
sinc functions of modulus at most one. Integrate on this ball,
use local convergence of \(j_b\), and increase the radius.

Write \(\lambda_{j,n}\) for the levels of the Rayleigh quotient
\(pD_n(v)/\langle v,Bv\rangle\). The sampled first \(j\) continuum
levels give \(\lambda_{j,n}\le\mu_j+O_{R,j}(b^{s/2})\). For the
lower bound take a subsequence realizing the lower limit and extract
strong limits of the first \(j\) vectors normalized in the
\(B\)-inner product. Their multiples by \(\sqrt p\) tend to an
orthonormal family. Applying the form lower bound to every linear
combination and then applying continuum min--max gives
\(\liminf\lambda_{j,n}\ge\mu_j\). The limits are eigenvectors at
their respective levels. Indeed the $i$th limit has energy at most
\(\mu_i\) and is orthogonal to the previously obtained eigenvectors,
so it attains the minimum on that complement. This argument allows
arbitrary choices within repeated eigenspaces.

We need a rate only for the first vector. Put \(v=\ell/\|\ell\|_2\)
and \(V_n=U_n v\). The upper comparison and the denominator estimate
give \(D_n(v)\le\mu_1+O_R(b^{s/2})\). Let \(r=b^{1/4}\) and
smooth \(V_n\) by a nonnegative smooth probability density of
support radius \(r\), obtained by rescaling a fixed density.
The result \(z\) is supported in the closure of \(\Omega_r=(-r,1+r)^d\) and
\(\|z-V_n\|_2=O_R(r^{s/2})\). On \(|\xi|\le b^{-1/2}\) the
cell multiplier and the Fourier multiplier of this convolution
each have modulus at most one. Thus this part of \(\mathcal D(z)\)
is at most \((1+O_R(b))D_n(v)\). On the complement, rapid
decay of the smooth multiplier gives \(O(b^2)\) using
\(\|V_n\|_2=1\). The ground level and spectral gap on \(\Omega_r\)
are \((1+2r)^{-s}\) times their values on \(\Omega\).
The Rayleigh excess of \(z/\|z\|_2\) above this ground level is
therefore \(O_R(r^{s/2}+r+b)\). As \(z\ge0\), the gap implies
that its distance to the normalized, dilated ground state \(w_r\)
is \(O_R((r^{s/2}+r+b)^{1/2})\). To bound
\(\|w_r-w\|_2\), smooth the compactly supported \(w\)
at radius \(r\) with error \(O(r^{s/2})\), and compare the dilates
of the smooth function by integrating their derivatives in the
dilation parameter. The cost is
\(O(r)(1+\|\nabla(w*\rho_r)\|_2)=O(r^{s/2})\) by the Fourier
form bound and bounded support. We obtain
\(\|V_n-w\|_2=O_R(b^{s/16})\). Normalization gives
\(\|\ell\|_2=p^{-1/2}+O_R(b^{s/2})\), and the estimate for
\((B-pI)\ell\) proves the profile assertion.

For the reduced covariance, put
\(e_n=\zeta\kappa_R/(pn^s)\) and \(r_n=(1-1/c')/e_n\).
The eigenvalues of \(E=B^{-1/2}AB^{-1/2}\) are
\(\eta_{j,n}=1-e_n\lambda_{j,n}>0\).
The pole displacement prescribed in \cite[Section~2]{base} is
\(c_{*,n}-c'=\theta_0\lambda^2\bar g_N(R/n)^s\), with fixed
\(\theta_0>0\) and \(\bar g_N=O_R((\log n)^{-1})\).
Hence \(r_n=\lambda_{1,n}+o_R(1)\). If \(\ell_{j,n}\) are the
\(B\)-orthonormal generalized eigenvectors, the spectral decomposition
of \(G=a_R B^{1/2}(E^{-1}-c'I)^{-1}B^{1/2}\) gives
\begin{equation}
 n^{-s}G_h=\frac{a_Rp}{\zeta\kappa_R}
 \sum_{j\ge2}\frac{\eta_{j,n}}{c'(\lambda_{j,n}-r_n)}
       (B\ell_{j,n})(B\ell_{j,n})^{\mathsf T}.
\end{equation}
The operator norm of the tail from level \(m\) is bounded by
\(C_R/(\lambda_{m,n}-r_n)\), since \(B^{1/2}\ell_{j,n}\) are
orthonormal and \(B\) is uniformly bounded. Its upper limit tends
to zero as \(m\to\infty\). In each fixed group of complete
eigenspaces the embedded columns tend subsequentially to
\(\sqrt p\) times an orthonormal continuum eigenbasis by the
min--max comparison and the denominator estimate. The resulting
spectral projections are independent of the chosen basis. Pairing
with the strongly convergent tests and then removing the truncation
gives the stated limit.
\end{proof}

\subsection{Sine coefficients of the fractional ground state}\label{app:sine-support}

\noindent We prove that $w$ has infinitely many nonzero sine
coefficients for $0<\alpha<1$. Suppose that $w$ is a finite sine sum.
Write $T_t$ for the whole-space heat semigroup applied to its zero
extension and restricted to $\Omega$, and $T^D_t$ for killed heat
flow in $\Omega$. The maximum principle gives $T_t w\ge T^D_t w$
for $w\ge0$. The scalar identity
$x^\alpha=c_\alpha\int_0^\infty(1-e^{-tx})t^{-1-\alpha}\,dt$
gives in the interior
\begin{equation}
 (-\Delta_D)^\alpha w-H_\alpha w
   =c_\alpha\int_0^\infty (T_t w-T^D_t w)t^{-1-\alpha}\,dt .
\end{equation}
One can first test this identity against nonnegative smooth functions
of compact support. The integral over $1\le t\le2$ is bounded below
near the centre of a face by a positive constant. Indeed $T_t w$
there has a positive lower bound by its Gaussian kernel, while the
finite sine sum $T_t^D w$ vanishes uniformly at the face.
On the left both the finite sine sum $(-\Delta_D)^\alpha w$ and
$H_\alpha w=\mu_1 w$ vanish there. This contradiction rules out a
finite sine expansion.

\section{Gaussian source bounds and terminal estimates}\label{app:slices}
\subsection{Covariance bounds for a smooth source}

To prove the norm bounds in Lemma~\ref{lem:sources}, let $b$ be the
length scale of a covariance slice and put
\begin{equation}
 \rho_b(x)=\min\{b,R+\operatorname{dist}(x,\Lambda_n^c)\},\qquad
 \|v\|_{b,\omega}=\max_{|\nu|\le q}\sup_{x\in\Lambda_n}
   (b/\rho_b(x))^\omega\rho_b(x)^{|\nu|}|\nabla^\nu v(x)|.
\end{equation}
Here $\nabla^\nu$ is a lattice difference of multi-order $\nu$ on a
stencil inside the box, and $q$ is the fixed derivative order in the
terminal field norm. Use the weight $\omega$ of the matched trajectory,
with $\max\{0,(s-1)/2\}<\omega<\alpha$ and $\omega\ne1/2$.
In dimension one this weight is $\omega=1/8$
\cite[Proposition~7.14]{base}. In every dimension the terminal norm
$E_n$ is the raw norm $F_{n,\omega}$ of \cite[Section~7.1]{base},
denoted here by $\|\cdot\|_{n,\omega}$. Its order-zero term controls
the supremum norm. Fix $\omega'$ with $\omega<\omega'<\alpha$.

Corollary~2.13 of \cite{base} treats the integrated uniform source.
For a signed smooth source we use its two covariance-leg estimates
and keep the differences of the test function in the interior.
First use the original split, with integrated covariance $C_<^0$.
A slice $C_j$ of range $O(b)$ has $O(b^d)$ sites per row. Applying
Cauchy--Schwarz to the two covariance legs, each of size
$O_R(b^{-s/2})$ at order zero, bounds the absolute row sum by
$C_R b^s$. With $k$ differences on the first leg, the slice comparison
and the weighted stencil estimates of
\cite[Lemma~2.8, Proposition~2.9 and Lemmas~2.10--2.11]{base}
give, for $f=f_n(F)$ with $\|f\|_\infty\le C_F n^{-s}$,
\begin{equation}
 |\nabla^k(C_j f)(x)|\le C_{R,F} n^{-s} b^s
       \rho_b(x)^{-k}(\rho_b(x)/b)^{\omega'},\qquad 0\le k\le q .
\end{equation}
Within distance $O(b)$ of a face, $\rho_n\lesssim\rho_b$. Multiplying
by the scale-$n$ stencil weight bounds the last display by
$C_{R,F}n^{-s+\omega}b^{s-\omega}$. In the interior the covariance
row agrees with its bulk convolution, so the differences act on $f$.
There $\|\nabla^k f\|_\infty\le C_F n^{-s-k}$ on the support of
the row and $b\lesssim\rho_n\le n$, giving the same bound. The
finitely many merged microscopic slices cost an $R$-dependent
constant. Summing geometrically, using $\omega<s$ and $b\lesssim n$,
gives $\|C_<^0 f\|_{E_n}=O_{R,F}(1)$. Summing the row bounds gives
$\|C_<^0\|_{2\to2}=O_R(n^s)$ by symmetry.

In dimension four use the unperturbed covariances $C_<^0,C_T^0$ at
cut $k$. The transfer identity gives
\begin{equation}
 C_<=tC_<^0+(1-t)a_R A+t\eta M,\qquad
 C_T=t(C_T^0-\eta M)
\end{equation}
by \cite[Lemma~5.4, Eq.~(5.19)]{base}, applied at this cut in
the proof of Corollary~7.13. Here
$\eta$ is bounded, $0\le M\le a_R A$, and $M\le C_R C_T^0$ for the
original terminal covariance $C_T^0$ at this cut. Both extra terms
have bounded operator norm. The finite-range kernel $A$ acts on $f$
by bulk convolution outside a microscopic collar. Inside it the
weighted differences cost $O_{R,F}(n^{\omega-s})$.
For any weighted stencil $D$ in $\|\cdot\|_{n,\omega}$, the
conditioned-field and profile bounds at the earlier cut in \cite[proof of Corollary~7.13]{base} give
\begin{equation}
 D C_T^0 D^{\mathsf T}\le C_R(n^{-2}+\sigma_0^2 n^{-4}),\qquad
 \sigma_0^2=\ell^{\mathsf T}C_T^0\ell=O_R(n^2\log n)
 .
\end{equation}
Thus $|DMf|\le(DMD^{\mathsf T})^{1/2}(f^{\mathsf T}Mf)^{1/2}
\le C_{R,F}(n^{-1}+\sigma_0n^{-2})$. This proves both estimates for
$C_<$ under the transfer.

Finally \eqref{eq:terminaldata} gives
$\E_{C_h}\|Y\|_{E_n}^2=O_R(n^{-s})$, hence
$\|C_h\|_{2\to2}\le\E_{C_h}\|Y\|_2^2=O_R(Vn^{-s})=O_R(n^s)$.
In $C_h f=\E_{C_h}[Y\ip{Y}{f}]$, Cauchy--Schwarz gives
$\|C_h f\|_{E_n}\le(\E_{C_h}\|Y\|_{E_n}^2)^{1/2}(f^{\mathsf T}C_h f)^{1/2}
=O_{R,F}(1)$. Adding the estimates for $C_<$ gives the bounds for
$K$ in \eqref{eq:cov-split}.

\subsection{Terminal estimates at the matched point}\label{app:terminal-inputs}

For $d=1,2,3$, Proposition~7.14 of \cite{base} verifies
Assumption~7.1 there at the physical temperature used in
Theorem~\ref{thm:tests}. It gives $g=g_N\asymp_R(\log V)^{-1}$,
the variance, profile and conditioned Gaussian moment in
\eqref{eq:terminaldata}, and the pointwise bound on the remainder
and its first two derivatives in arbitrary field directions.
In particular, the terminal packing estimate
\cite[Proposition~7.2]{base} gives
\begin{equation}
 \epsilon_n\le C_R\bigl(\sqrt g+g^{-9/4}\|K_N^{\rm RG}\|\bigr)
 \le C_R\bigl(\sqrt g+g^{3/4}+U d_0^N g^{-9/4}\bigr)=o(1).
\end{equation}
Here $K_N^{\rm RG}$ is the terminal polymer activity, $U=O_R(1)$
is the initial activity bound, $0<d_0<1$ and $N\asymp\log n$.
Integrating the pointwise bound over the conditioned field $Y$, as
in the proof of \cite[Proposition~7.4]{base}, gives
\eqref{eq:envelope} with the same factors $\|v_i\|_{E_n}/h_*$.

For $d=4$, use the temperature and change of field in
\cite[Corollary~7.13]{base}. Its proof uses the prescribed scalar
density at cut $J$ to select the mass. At the earlier cut $k$, the
terminal coefficients of the transformed field satisfy the bounds of
\cite[Proposition~7.2]{base} at that mass, with
$g=g_k\asymp_R(\log n)^{-1}$. Lemma~5.4 of \cite{base}, applied
at the fixed change of aspect as in the proof of Corollary~7.13,
gives $\sigma_T^2\asymp_R n^2/g_k$, the profile bound for $u$ and
the exponential moment of $Y$ in \eqref{eq:terminaldata}. The
same integration of the two-derivative bound gives \eqref{eq:envelope}
for the field in \eqref{eq:cov-split}. Thus the moment calculation
uses only the coefficient estimates at the selected mass.

\section{Displacement of the pseudocritical point}\label{app:shift}

Recall that $\tau_c=\tau_c(d,R)=\inf\{\tau>0:\langle\sigma_0\rangle_\tau^+>0\}$
and $\Delta_n=\tau_{\Lambda_n}^{\rm pc}-\tau_c$ and $(u)_+=\max\{u,0\}$.

\begin{proof}[Proof of Theorem~\ref{thm:pc-shift}]
Let $G_{n,\tau}^{\rm F}$ be the zero-field free-box two-point matrix at
physical inverse temperature $\tau$ and put
$\chi_n^{\rm unif}(\tau)=V^{-1}\one^{\mathsf T}G_{n,\tau}^{\rm F}\one$.
At the matched point $G_{n,\tau_{\Lambda_n}^{\rm pc}}^{\rm F}=G_{\rm sp}$ and
\cite[Theorem~1.1]{base} gives
\begin{equation}\label{eq:shift-matched}
 \chi_n^{\rm unif}(\tau_{\Lambda_n}^{\rm pc})\asymp_{d,R}n^s\sqrt{\log n}.
\end{equation}
The ferromagnetic pair kernel is
$K_R(z)=-(L_R^\alpha)(0,z)$ for $z\ne0$, with $K_R(0)=0$.
Thus $K_R(x-y)=-(J_n)_{xy}$ at distinct sites in the box. When
$\alpha<1$ the binomial series gives
\begin{equation}\label{eq:shift-walk}
 I-L_R^\alpha=\sum_{k\ge1}a_kD_R^{*k},\qquad
 a_k=(-1)^{k+1}\binom{\alpha}{k}>0,\qquad \sum_{k\ge1}a_k=1.
\end{equation}
At $\alpha=1$, $K_R=D_R$. Hence $j_R=\sum_zK_R(z)$ lies in $(0,1]$.
For $\tau j_R<1$ the Dobrushin influence sum is at most
$\sum_z\tanh(\tau K_R(z))\le\tau j_R<1$. Uniqueness and spin-flip
symmetry imply $\tau_c\ge j_R^{-1}\ge1$. The matching construction in
\cite[Corollaries~6.3 and~7.13, Proposition~7.14]{base} takes place in
$|\beta-1|\le A_dR^{-d}$, with $A_d$ fixed before $R$ and $n$.
Since $\tau=\zeta\beta$ and $0<\zeta\le1$ in Section~1, this gives
$\tau_{\Lambda_n}^{\rm pc}\le1+A_dR^{-d}$ and hence the stated
$R^{-d}$ bound in all four dimensions.

Set $h_x(\tau)=\sum_{y\in\Lambda_n}G_{n,\tau}^{\rm F}(x,y)$ and
$H_n(\tau)=\max_{x\in\Lambda_n}h_x(\tau)$. The bulk input is the
free-state estimate
$\langle\sigma_0\sigma_z\rangle_{\tau_c}^{\rm f}
\le C_{d,R}(1+|z|_\infty)^{-s}$ of \cite[Theorem~1]{park-bulk}.
By Griffiths' domain monotonicity the free-box correlations at $\tau_c$
satisfy the same bound. Summing over shells about $x$ gives
\begin{align}
 H_n(\tau_c)\le C_{d,R}\sum_{k=0}^{n-1}(1+k)^{d-1-s}\le A_{d,R}n^s,
\end{align}
since $d-s=s$. At the matched point the maximum row sum dominates its
average, so \eqref{eq:shift-matched} gives
$H_n(\tau_{\Lambda_n}^{\rm pc})\ge a_{d,R}n^s\sqrt{\log n}$.
For large $n$ this exceeds $H_n(\tau_c)$. Monotonicity in $\tau$
then gives $\tau_{\Lambda_n}^{\rm pc}>\tau_c$.

For brevity write $G=G_{n,\tau}^{\rm F}$. Differentiating the spin
expectation with respect to $\tau$ and applying Lebowitz's zero-field
four-point inequality \cite{lebowitz} gives
\begin{align}
 \frac{d}{d\tau}G(x,y)
 =\sum_{\{u,v\}\subset\Lambda_n} K_R(u-v)
      \operatorname{Cov}_{n,\tau}^{\rm F}(\sigma_x\sigma_y,\sigma_u\sigma_v)
 \le \sum_{u,v\in\Lambda_n} G(x,u)K_R(u-v)G(v,y).
\end{align}
The two cross-pairings turn the unordered bond sum into the ordered
sum on the right. Summing over $y$ yields
$h_x'(\tau)\le j_R h_x(\tau)H_n(\tau)$. The maximum of these finitely
many smooth row sums is locally Lipschitz and satisfies
$H_n'\le j_R H_n^2$ almost everywhere. Integrating the differential
inequality for $1/H_n$ gives
\begin{align}
 \Delta_n\ge \frac1{j_R}
 \left(\frac1{H_n(\tau_c)}-\frac1{H_n(\tau_{\Lambda_n}^{\rm pc})}\right)
 \ge \frac{n^{-s}}{j_R}
 \left(\frac1{A_{d,R}}-\frac1{a_{d,R}\sqrt{\log n}}\right).
\end{align}
The last parenthesis is bounded below by a positive constant for
sufficiently large $n$.

For \eqref{eq:pc-shift-upper} take $d<4$. If $\Delta_n\le0$ the bound
is immediate. Otherwise set $\tau=\tau_{\Lambda_n}^{\rm pc}$,
$\theta=(\tau_c+\tau)/2$, and let $D_r$ be the uniform kernel of
Section~1 with integer spread $1\le r\le n$.
Two estimates on $K_R$ will be useful. For a nonnegative symmetric
coupling $H$ with zero diagonal, denote its missing bond weight per site by
$e_n(H)=V^{-1}\sum_{x\in\Lambda_n,\,y\notin\Lambda_n}H(x-y)$.
At fixed $R$,
\begin{align}
 \min_{0<|z|_\infty\le r}K_R(z)\ge c_{d,R}r^{-d-s},
 \qquad e_n(K_R)\le C_{d,R}
 \begin{cases}
 n^{-s},&s<1,\\
 n^{-1}\log(2n),&s=1,\\
 n^{-1},&s>1 .
 \end{cases}
\end{align}
For the first estimate, $a_k\asymp_d k^{-1-s/2}$ in \eqref{eq:shift-walk},
and the local central limit theorem for the aperiodic spread-out walk
gives $D_R^{*k}(z)\ge c_{d,R}k^{-d/2}$ for $r^2\le k\le2r^2$ and
$|z|_\infty\le r$ and large $r$. Positivity in \eqref{eq:shift-walk}
covers the finitely many smaller $r$. For the second estimate, a displacement $z$ leaves
the cube from at most a fraction $\min\{1,d|z|_\infty/n\}$ of its sites.
The walk after $k$ steps has expected norm at most $C_{d,R}\sqrt{k}$.
Splitting $\sum_{k\ge1}a_k\min\{1,C_{d,R}\sqrt{k}/n\}$ at $k=n^2$
gives the three bounds on $e_n(K_R)$.

Choose $\delta=c_{d,R}(\tau-\theta)r^{-s}$ with the constant small
enough that $H_0=\theta K_R$ and $H_1=H_0+\delta D_r$ satisfy
$H_0\le H_1\le\tau K_R$ entrywise. Let $P_n(H)$ and $P_\infty(H)$
be the free-box and bulk log partition functions per site with pair
coupling $H$. Tiling a large cube by $n$-boxes gives
$0\le P_\infty(H)-P_n(H)\le e_n(H)/2$, since adding the bonds between
tiles increases the pressure by at most their total weight per site.
Interpolation between $H_0$ and $H_1$, Griffiths monotonicity and
$\|D_r\|_\infty\le C_dr^{-d}$ give
\begin{align}
 P_n(H_1)-P_n(H_0)
 \le \frac{\delta}{2V}\sum_{x,y\in\Lambda_n}
                D_r(x-y)G_{n,\tau}^{\rm F}(x,y)
 \le C_d\delta r^{-d}\chi_n^{\rm unif}(\tau).
\end{align}
Write $m_+(\theta)=\langle\sigma_0\rangle_\theta^+$. In the bulk,
$P_\infty(H_1)-P_\infty(H_0)\ge\delta m_+(\theta)^2/2$.
Indeed, keep $n,r$ fixed and use the same $H_0$ plus
boundary field in both partition functions on an auxiliary $M$-box.
Positive association and plus-volume monotonicity bound every
$H_0$ pair expectation below by the squared plus magnetization.
Jensen's inequality for the added interior bonds gives the claim as
$M\to\infty$, since $D_r$ has total mass one and the boundary weight
per site tends to zero by summability. Combining the pressure estimates
and using $e_n(H_1)\le\tau e_n(K_R)$ yields
\begin{align}
 0<m_+(\theta)^2
 \le C_d r^{-d}\chi_n^{\rm unif}(\tau)+\tau e_n(K_R)/\delta .
\end{align}
By the magnetization inequality \cite[Theorem~2.1(1)]{sharpness},
\begin{align}
 m_+(\theta)^2\ge 1-(\tau_c/\theta)^2
   =\frac{\Delta_n(\theta+\tau_c)}{2\theta^2}
   \ge\frac{\Delta_n}{2\theta}.
\end{align}
Here $\theta\le\tau\le 1+A_dR^{-d}$ and
$\delta=c_{d,R}\Delta_n r^{-s}/2$. Multiplying the resulting
inequality by $\Delta_n$ and taking the positive quadratic root gives
\begin{align}
 (\Delta_n)_+\le C_{d,R}\left(
 \frac{\chi_n^{\rm unif}(\tau_{\Lambda_n}^{\rm pc})}{r^d}
     +\sqrt{e_n(K_R)r^s}\right).
\end{align}
Using \eqref{eq:shift-matched}, choose $r$ to balance the two terms.
Up to integer rounding, the choices are $n^{3/5}(\log n)^{2/5}$ for
$d=1$, $n^{3/5}$ for $d=2$, and $n^{8/15}(\log n)^{2/15}$ for $d=3$.
Substitution gives \eqref{eq:pc-shift-upper}.

To conclude convergence, for each fixed $0<\tau<\tau_c$,
subcritical sharpness \cite[Theorem~2.1(2)]{sharpness} and domain
monotonicity bound $\chi_n^{\rm unif}(\tau)$ by the finite bulk
susceptibility. By \eqref{eq:shift-matched} and monotonicity,
$\tau_{\Lambda_n}^{\rm pc}>\tau$ eventually. This gives
$\liminf_n\tau_{\Lambda_n}^{\rm pc}\ge\tau_c$ in every dimension.
For $d<4$, \eqref{eq:pc-shift-upper} supplies the corresponding limsup.

For $d=4$, $K_R=D_R$ has finite range and contains every nearest-neighbour
bond. At any fixed $\tau>\tau_c$, slab percolation and equality of the
free and wired bulk FK measures hold \cite[Sections~2.1--2.3]{bodineau-finite}.
These are the hypotheses of the finite-range FK--Ising crossing-cluster
estimate \cite[Section~3.2, Eq.~(3.6)]{crossing-cluster}.
Take a central block of dyadic side between $n/8$ and $n/4$. For large
$n$ all bonds touching the block lie in $\Lambda_n$. The estimate is
uniform over the boundary wiring induced by bonds elsewhere in the
free cube. Thus, with probability at least $1/2$, a cluster of the
free cube contains at least $b_{\tau,R}V$ sites for some $b_{\tau,R}>0$.
The Edwards--Sokal identity gives
$\chi_n^{\rm unif}(\tau)=V^{-1}\E_{\rm FK}\sum_{\mathcal C}|\mathcal C|^2
\ge b_{\tau,R}^2V/2$, the sum running over the clusters of the free cube.
Since \eqref{eq:shift-matched} gives
$\chi_n^{\rm unif}(\tau_{\Lambda_n}^{\rm pc})=O_R(n^2\sqrt{\log n})=o(V)$,
monotonicity puts $\tau_{\Lambda_n}^{\rm pc}<\tau$ eventually for every
fixed $\tau>\tau_c$. This proves convergence in dimension four.
\end{proof}

\section*{AI Statement}

This document is generated with a collaboration with ChatGPT-Astra 6.0 and ChatGPT-Astra 6.2.


\begin{thebibliography}{99}
\bibitem{base}
J.~Park, \emph{Pseudocritical point of the spread-out Ising model with free
boundaries at the upper critical dimension in $1\le d\le4$},
GitHub preprint,
\href{https://github.com/jiwoonGit/phi4/blob/main/FSS/FSS_FBC-dim-extension.pdf}{DOI:10.5281/zenodo.24167641} (2026).

\bibitem{flores}
E.~Flores-Sola, B.~Berche, R.~Kenna and M.~Weigel,
\href{https://doi.org/10.1103/PhysRevLett.116.115701}
{\emph{Role of Fourier modes in finite-size scaling above the upper
critical dimension}}, Phys. Rev. Lett. \textbf{116} (2016), 115701.

\bibitem{fisher-barber}
M.~E.~Fisher and M.~N.~Barber,
\href{https://doi.org/10.1103/PhysRevLett.28.1516}
{\emph{Scaling theory for finite-size effects in the critical region}},
Phys. Rev. Lett. \textbf{28} (1972), 1516--1519.

\bibitem{brezin-zinnjustin}
E.~Br\'ezin and J.~Zinn-Justin,
\href{https://doi.org/10.1016/0550-3213(85)90379-7}
{\emph{Finite size effects in phase transitions}},
Nucl. Phys. B \textbf{257} (1985), 867--893.

\bibitem{rudnick}
J.~Rudnick, G.~Gaspari and V.~Privman,
\href{https://doi.org/10.1103/PhysRevB.32.7594}
{\emph{Effect of boundary conditions on the critical behavior of a
finite high-dimensional Ising model}}, Phys. Rev. B \textbf{32}
(1985), 7594--7597.

\bibitem{wittmann}
M.~Wittmann and A.~P.~Young,
\href{https://doi.org/10.1103/PhysRevE.90.062137}
{\emph{Finite-size scaling above the upper critical dimension}},
Phys. Rev. E \textbf{90} (2014), 062137.

\bibitem{honchar}
Yu.~Honchar, B.~Berche, Yu.~Holovatch and R.~Kenna,
\href{https://doi.org/10.5488/CMP.27.13603}
{\emph{When correlations exceed system size: finite-size scaling in
free boundary conditions above the upper critical dimension}},
Condens. Matter Phys. \textbf{27} (2024), 13603.

\bibitem{luijten}
E.~Luijten and H.~W.~J.~Bl\"ote,
\href{https://doi.org/10.1103/PhysRevB.56.8945}
{\emph{Classical critical behavior of spin models with long-range
interactions}}, Phys. Rev. B \textbf{56} (1997), 8945--8958.

\bibitem{lv}
J.-P.~Lv, W.~Xu, Y.~Sun, K.~Chen and Y.~Deng,
\href{https://doi.org/10.1093/nsr/nwaa212}
{\emph{Finite-size scaling of $O(n)$ systems at the upper critical
dimensionality}}, Natl. Sci. Rev. \textbf{8} (2021), nwaa212.

\bibitem{park-torus}
J.~Park, \emph{Torus scaling limits and the plateau of the critical
weakly coupled $|\varphi|^4$ model in $d\ge4$},
\href{https://arxiv.org/abs/2511.06321}{arXiv:2511.06321} (2025).

\bibitem{camia-jiang-newman}
F.~Camia, J.~Jiang and C.~M.~Newman,
\href{https://doi.org/10.1007/s00440-021-01041-9}
{\emph{The effect of free boundary conditions on the Ising model in high
dimensions}}, Probab. Theory Related Fields \textbf{181} (2021), 311--328.

\bibitem{duminil-panis}
H.~Duminil-Copin and R.~Panis,
\href{https://doi.org/10.1007/s00220-025-05236-2}
{\emph{New lower bounds for the (near) critical Ising and $\varphi^4$
models' two-point functions}}, Comm. Math. Phys. \textbf{406}
(2025), article~56.

\bibitem{michta-park-slade}
E.~Michta, J.~Park and G.~Slade,
\href{https://doi.org/10.1002/cpa.22256}
{\emph{Boundary conditions and universal finite-size scaling for the
hierarchical $|\varphi|^4$ model in dimensions 4 and higher}},
Comm. Pure Appl. Math. \textbf{78} (2025), 2001--2118.

\bibitem{park-slade}
J.~Park and G.~Slade,
\href{https://doi.org/10.1007/s00023-025-01566-y}
{\emph{Boundary conditions and the two-point function plateau for the
hierarchical $|\varphi|^4$ model in dimensions 4 and higher}},
Ann. Henri Poincar\'e \textbf{27} (2026), 1497--1561.

\bibitem{liu-park-slade}
Y.~Liu, J.~Park and G.~Slade,
\href{https://doi.org/10.1103/xq7j-zll9}
{\emph{Universal finite-size scaling in high-dimensional critical
phenomena}}, Phys. Rev. E \textbf{113} (2026), 024143,
\href{https://arxiv.org/abs/2412.08814}{arXiv:2412.08814}.

\bibitem{park-bulk}
J.~Park, \emph{Correlation function of the fractional spread-out Ising
models at their upper critical dimension and critical temperature},
preprint (2026), Theorem~1.

\bibitem{lebowitz}
J.~L.~Lebowitz,
\href{https://cmsr.rutgers.edu/images/people/lebowitz_joel/publications/jll.pub_101.pdf}
{\emph{GHS and other inequalities}}, Comm. Math. Phys. \textbf{35}
(1974), 87--92.

\bibitem{sharpness}
H.~Duminil-Copin and V.~Tassion,
\href{https://arxiv.org/abs/1502.03050}{\emph{A new proof of the sharpness
of the phase transition for Bernoulli percolation and the Ising model}},
Comm. Math. Phys. \textbf{343} (2016), 725--745.

\bibitem{bodineau-finite}
T.~Bodineau, \href{https://arxiv.org/abs/math/0409079}{\emph{Translation
invariant Gibbs states for the Ising model}}, Probab. Theory Related
Fields \textbf{135} (2006), 153--168.

\bibitem{crossing-cluster}
T.~Bodineau, D.~Ioffe and Y.~Velenik,
\href{https://www.unige.ch/~velenik/papers/BIV01.pdf}{\emph{Winterbottom
construction for finite range ferromagnetic models: an $L_1$-approach}},
J. Stat. Phys. \textbf{105} (2001), 93--131.
\end{thebibliography}
\end{document}
